% Template for one-page artifact description appendices in lncs volumes
% Written by Camil Demetrescu and Erik Ernst
% April 8, 2014

% ARTIFACT: do not change the following line
\clearpage\section{Artifact Description}

% ARTIFACT: please note that the rest of this file contains (1) LaTeX commands
% and environments that you must use in the way they are given here, and in the
% same order, each producing one section of text; (2) example data filling in
% those commands and environments with text to show how they may be used; you
% must preserve the former and replace the latter by information concerning
% your artifact; the comments describe each of these sections to indicate
% their purpose; also note that the entire appendix must fit on one page

% ARTIFACT: section on authors of the artifact; note that this may differ from
% the list of authors of the paper, e.g., you may wish to mention other people
% who contributed to the artifact, but who did not co-author the paper
\artifactauthors{Design and documentation: Author One, Author Two,
  and Author Three. Core developer: Author Three.}

% ARTIFACT: section giving a brief summary of the purpose of the artifact
\begin{summary}
  The artifact is based on {\tt ArtiFact}, a dynamic program analysis tool that
  can profile paths spanning multiple loop iterations in the control flow
  graph of a Java program. Profiled paths are obtained as the concatenation of
  up to k acyclic paths, where k is a user-defined parameter. The profiler was
  implemented as a patch to the Jikes Research Virtual Machine. The provided
  package is designed to support repeatability of the experiments of the
  paper: in particular, it allows users to test the profiler on a variety of
  benchmarks and includes detailed instructions and scripts for running them
  and for visualizing and interpreting collected profiles. Instructions for
  rebuilding the profiler from scratch in the Jikes RMV are also provided.
\end{summary}

% ARTIFACT: section on the contents of the artifact (code, data, etc.)
\begin{content}
  The artifact package includes:
  \begin{itemize}
  \item a Jikes RVM patch that allows users to profile hot methods of Java programs;
  \item a Java program for visualizing the reports generated by the profiler using GraphViz;
  \item detailed instructions for using the artifact and for rebuilding it from scratch, provided as an {\tt index.html} file.
  \end{itemize}
  To simplify repeatabiliy of our experiments, we provide a VirtualBox disk
  image containing a lightweight Linux distribution fully configured for
  testing our profiler. The image contains Precise Puppy 5.6 - a stripped-down
  Linux distribution based on Ubuntu 12.04 LTS (Precise Pangolin). In
  particular, since the process of building the Jikes RVM can be
  time-consuming (several tools, such as the GNU Classpath and JUnit, are
  downloaded and built during the first source compilation), this image comes
  up with an already patched and compiled Jikes RVM build, and includes all
  benchmarks used in the paper. The Jikes RVM has been compiled with OpenJDK
  1.6.0\_27 and gcc 4.6.3 using the default production configuration as
  suggested on the project's website. 
\end{content} 

% ARTIFACT: section containing links to sites (e.g., github) holding the
% latest version of the code, if any
\begin{getting}
  The artifact endorsed by the Artifact Evaluation Committee is available free
  of charge as supplementary material of this paper on SpringerLink. The
  latest version of our code is available on the Jikes RVM Research Archive:
  {\url http://sourceforge.net/p/jikesrvm/research-archive/41/} and at 
  {\url http://www.cs.fu.utopia/$\sim$authorthree/projects.html}.
\end{getting} 

% ARTIFACT: section specifying the platforms on which the artifact is known to
% work, including requirements beyond the operating system such as large
% amounts of memory or many processor cores
\begin{platforms}
  The artifact is known to work on any platform running Oracle VirtualBox
  version~4 ({\url https://www.virtualbox.org/}) with at least 5~GB or free
  space on disk and at least 2~GB of free space in RAM.
\end{platforms}

% ARTIFACT: section specifying the license under which the artifact is
% made available
\license{EPL-1.0 ({\url http://www.eclipse.org/legal/epl-v10.html})}

% ARTIFACT: section specifying the md5 sum of the artifact master file
% uploaded to SpringerLink, enabling downloaders to check that the file is the
% expected version and suffered no damage during download
\mdsum{3fda3a372cea21a28d433519a4412ded}

% ARTIFACT: section specifying the size of the artifact master file uploaded
% to SpringerLink
\artifactsize{1.2 GB}

% ARTIFACT: include here any additional references, if needed...

%\bibliographystyle{unsrt}
%\begin{thebibliography}{10}
%...
%\end{thebibliography}

%%% Local Variables: 
%%% mode: latex
%%% TeX-master: "main"
%%% End: 

